\chapter{Conclusions}\label{ch:conclusions}
\section{Summary of work}
RayGuard has a concrete single-scan prototype path and a growing research
software foundation. Audits, strict contracts and bounded real execution
support image-to-result integration. The separate application exposes localized
findings and traceable operator review. Evaluation and initialization software
have been integrated. The reduced generic study completed six attempted
configurations, two 100-epoch confirmations and a frozen full published-test
comparison, with saved predictions and independently recomputed metrics.

\section{Critical appraisal}
The current evidence includes reproducible software checks and measured generic
detector performance, with a substantial validation-to-test AP gap. The dataset audits reveal why broad accuracy
or unseen-data claims require care: task exports disagree, exact images cross
the supplied split, boundary findings remain, and physical groups are not
established locally. The collection inventory is useful progress evidence,
but it does not yet establish independent labeled cases or accepted completion.

The project has not demonstrated full P1 host/threat decisions, distributed-case
fusion, a live scanner connection or repeatable benchmark improvement. The
preselected candidate achieved higher test AP in one run, but lower small-object
AP and fixed-threshold recall; both models missed roughly half the annotated
boxes at confidence 0.25. These
limits affect the system's permitted interpretation. An operator should not
treat an empty detector output as clearance, nor a plausible component grouping
as proof of a functional device. Explicit states and provenance help communicate
these limits but cannot replace performance and safety validation.

\section{Further work}
The immediate handoff is review of the qualified implementation, frozen results
and report evidence. No additional training is needed to finish the authorized
deadline scope. Future experiments should establish physical groups and assess
repeatability on training-derived validation before any new final test protocol.
The deferred specific-task and original twenty-attempt studies must not be
reported as completed. The report deadline remains 24 November 2026 at
17:00 UTC+4; the completed generic results and their limitations support the
23 November evidence freeze.

The longer-term distributed work requires QC, accepted collection counts,
component annotations, physical/case provenance and a case-disjoint evaluation
design. A transparent independent-scan aggregation baseline should be assessed
before introducing learned fusion, provided suitable ground truth exists.
Human acceptance on the presentation machine and an observed scanner handoff
remain separate deployment steps. The overall collection deadline of
30 November is retained independently of report submission.

The completed report should explain both the achieved prototype and the work
still necessary to justify stronger claims. Its final contribution table,
forms and conclusions require student and supervisor review.
